% =============================================================================
% Section 9: Evaluation
% =============================================================================
\section{Evaluation}
\label{sec:evaluation}

% ---------------------------------------------------------------------------
% 9.1 Quality Targets
% ---------------------------------------------------------------------------
\subsection{Quality Targets}
\label{sec:eval:quality}

The pipeline defines quantitative quality targets at each phase, with minimum
thresholds that block further progress and aspirational targets for optimal
performance.

\begin{table}[H]
  \centering
  \caption{Phase 1 quality targets (encoder adaptation).}
  \label{tab:quality_phase1}
  \begin{tabular}{lccc}
    \toprule
    \textbf{Metric} & \textbf{Target} & \textbf{Minimum} & \textbf{Status} \\
    \midrule
    Dice (WT) on \texttt{lora\_val} & $\geq 0.85$ & $\geq 0.80$ & \textit{[to fill]} \\
    Dice improvement over frozen BSF & $\geq 0.05$ & $\geq 0.02$ & \textit{[to fill]} \\
    Linear probe Vol $R^2$ (adapted) & $\geq 0.50$ & $\geq 0.30$ & \textit{[to fill]} \\
    \bottomrule
  \end{tabular}
\end{table}

\begin{table}[H]
  \centering
  \caption{Phase 2 quality targets (SDP).}
  \label{tab:quality_phase2}
  \begin{tabular}{lccc}
    \toprule
    \textbf{Metric} & \textbf{Target} & \textbf{Minimum} & \textbf{Status} \\
    \midrule
    Vol $R^2$ on validation & $\geq 0.90$ & $\geq 0.80$ & \textit{[to fill]} \\
    Loc $R^2$ on validation & $\geq 0.95$ & $\geq 0.85$ & \textit{[to fill]} \\
    Shape $R^2$ on validation & $\geq 0.40$ & $\geq 0.25$ & \textit{[to fill]} \\
    Max cross-partition corr & $< 0.20$ & $< 0.30$ & \textit{[to fill]} \\
    Dims with var $> 0.5$ & $\geq 95\%$ & $\geq 85\%$ & \textit{[to fill]} \\
    $\dCor(\text{vol}, \text{loc})$ & $< 0.10$ & $< 0.20$ & \textit{[to fill]} \\
    \bottomrule
  \end{tabular}
\end{table}

\begin{table}[H]
  \centering
  \caption{Phase 4 quality targets (growth prediction).}
  \label{tab:quality_phase4}
  \begin{tabular}{lccc}
    \toprule
    \textbf{Metric} & \textbf{Target} & \textbf{Minimum} & \textbf{Status} \\
    \midrule
    Vol prediction $R^2$ (LOPO, best model) & $\geq 0.70$ & $\geq 0.50$ & \textit{[to fill]} \\
    Calibration (95\% CI coverage, GP models) & 0.90--0.98 & $\geq 0.80$ & \textit{[to fill]} \\
    Per-patient trajectory $r$ ($n_i \geq 3$) & $\geq 0.80$ & $\geq 0.60$ & \textit{[to fill]} \\
    PA-MOGP $R^2$ $>$ H-GP $R^2$ & $> 0$ & --- & \textit{[to fill]} \\
    \bottomrule
  \end{tabular}
\end{table}

% ---------------------------------------------------------------------------
% 9.2 Ablation Matrix
% ---------------------------------------------------------------------------
\subsection{Ablation Study Design}
\label{sec:eval:ablations}

Ten ablation experiments systematically evaluate design decisions across all
pipeline phases:

\begin{table}[H]
  \centering
  \caption{Full ablation matrix (A1--A10).}
  \label{tab:ablation_matrix}
  \begin{tabular}{clcp{5.5cm}}
    \toprule
    \textbf{ID} & \textbf{Variable} & \textbf{Conditions} & \textbf{Primary metric} \\
    \midrule
    A1 & LoRA rank & $\{2, 4, 8, 16, 32\}$ & Dice, probe $R^2$ \\
    A2 & LoRA vs.\ DoRA & $\{\LoRA, \text{DoRA}\}$ at $r=8$ & Dice, probe $R^2$ \\
    A3 & Aux semantic heads & $\{$with, without$\}$ & Phase~2 $R^2$ \\
    A4 & SDP dimension & $\{64, 128, 256\}$ & $R^2$, dCor \\
    A5 & VICReg + dCor & $\{$full, no cov, no dCor, no both$\}$ & Cross-partition corr \\
    A6 & Growth model & $\{\LME, \HGP, \PAMOGP\}$ & Vol $R^2$ (LOPO) \\
    A7 & ComBat effect & $\{$with, without$\}$ & Vol $R^2$ (LOPO) \\
    A8 & GP mean function & $\{$zero, linear, Gompertz$\}$ & Vol $R^2$ (LOPO) \\
    A9 & GP kernel & $\{$M-3/2, M-5/2, SE$\}$ & Vol $R^2$ (LOPO) \\
    A10 & Cross-partition coupling & $\{$with, without$\}$ & Vol $R^2$, calibration \\
    \bottomrule
  \end{tabular}
\end{table}

\subsubsection{A6: Model Comparison}

The head-to-head comparison of the three GP models is the central evaluation
of the pipeline. The comparison uses LOPO-CV and reports:

\begin{description}[leftmargin=*]
  \item[Vol $R^2$:] Coefficient of determination between predicted and
    observed whole-tumour volume across all held-out observations.
  \item[Vol MAE:] Mean absolute error in physical volume (mm$^3$).
  \item[Latent MSE:] Mean squared error in the latent volume partition
    $\Vol$ (before decoding).
  \item[Calibration:] Fraction of true observations within the predicted
    95\% confidence interval (target: 0.90--0.98).
  \item[Per-patient $r$:] Pearson correlation between predicted and observed
    volume trajectories for patients with $n_i \geq 3$.
\end{description}

\begin{table}[H]
  \centering
  \caption{Model comparison results (LOPO-CV, 33 folds).}
  \label{tab:model_comparison}
  \begin{tabular}{lccccc}
    \toprule
    \textbf{Model} & \textbf{Params} & \textbf{Vol $R^2$} &
      \textbf{Vol MAE} & \textbf{Calibration} & \textbf{Per-patient $r$} \\
    \midrule
    LME     & 144 & \textit{[to fill]} & \textit{[to fill]} & --- & \textit{[to fill]} \\
    H-GP    & 72  & \textit{[to fill]} & \textit{[to fill]} & \textit{[to fill]} & \textit{[to fill]} \\
    PA-MOGP & 98  & \textit{[to fill]} & \textit{[to fill]} & \textit{[to fill]} & \textit{[to fill]} \\
    \bottomrule
  \end{tabular}
\end{table}

% ---------------------------------------------------------------------------
% 9.3 Figures
% ---------------------------------------------------------------------------
\subsection{Publication Figures}
\label{sec:eval:figures}

The evaluation produces 13 publication-quality figures documenting each
pipeline phase. Key figures include:

\begin{figure}[H]
  \centering
  \fbox{\parbox{0.8\textwidth}{\centering\vspace{3cm}
    \textbf{Figure placeholder:} Volume prediction scatter.\\
    Predicted vs.\ actual $\Delta V$ for all LOPO-CV test pairs,\\
    coloured by model (LME: blue, H-GP: orange, PA-MOGP: green).
  \vspace{3cm}}}
  \caption{Predicted vs.\ actual volume change ($\Delta V$, mm$^3$) across
    all LOPO-CV held-out observation pairs, coloured by model. The
    diagonal represents perfect prediction. Dispersion around the diagonal
    quantifies prediction error.}
  \label{fig:vol_scatter}
\end{figure}

\begin{figure}[H]
  \centering
  \fbox{\parbox{0.8\textwidth}{\centering\vspace{3cm}
    \textbf{Figure placeholder:} Trajectory prediction with uncertainty.\\
    3--4 example patients showing predicted trajectory\\
    (mean $\pm$ 95\% CI) overlaid on actual observations.
  \vspace{3cm}}}
  \caption{Growth trajectory predictions for selected patients. Solid lines
    show GP posterior mean; shaded regions show 95\% confidence intervals.
    Dots are actual observations. The GP uncertainty naturally widens with
    extrapolation distance.}
  \label{fig:trajectory_predictions}
\end{figure}

\begin{figure}[H]
  \centering
  \fbox{\parbox{0.8\textwidth}{\centering\vspace{3cm}
    \textbf{Figure placeholder:} Model comparison bar chart.\\
    $R^2$, MAE, and calibration for each model.
  \vspace{3cm}}}
  \caption{Head-to-head model comparison across three metrics. Error bars
    represent bootstrap 95\% confidence intervals computed over the 33
    LOPO-CV folds.}
  \label{fig:model_comparison}
\end{figure}

\begin{figure}[H]
  \centering
  \fbox{\parbox{0.8\textwidth}{\centering\vspace{3cm}
    \textbf{Figure placeholder:} Disentanglement matrix.\\
    Cross-partition correlation heatmap from SDP.
  \vspace{3cm}}}
  \caption{Cross-partition correlation heatmap after SDP training. Each cell
    shows the maximum absolute Pearson correlation between dimensions in the
    corresponding partition pair. Diagonal blocks (within-partition) show
    internal structure; off-diagonal blocks should approach zero for
    successful disentanglement.}
  \label{fig:disentanglement_matrix}
\end{figure}

\begin{figure}[H]
  \centering
  \fbox{\parbox{0.8\textwidth}{\centering\vspace{3cm}
    \textbf{Figure placeholder:} Learned GP length-scales per partition.\\
    Revealing characteristic timescales of volume/location/shape dynamics.
  \vspace{3cm}}}
  \caption{Learned GP kernel length-scales ($\ell$) for the PA-MOGP model.
    Volume dynamics have the longest characteristic timescale (slow growth),
    while shape dynamics are faster (shorter $\ell$). Location length-scales
    are expected to be very long (near-static centroid).}
  \label{fig:lengthscales}
\end{figure}

% ---------------------------------------------------------------------------
% 9.4 Methodology Revision (R1)
% ---------------------------------------------------------------------------
\subsection{R1 Methodology Revision}
\label{sec:eval:revision}

Empirical findings from Phases~1--2 motivated three principled revisions to
the pipeline architecture, collectively designated the R1 revision.

\subsubsection{Shape Partition Failure}

Across all experimental conditions, the shape partition achieved
$R^2 \leq 0.11$ for linear probes and approximately zero for RBF probes.
Cross-domain shape transfer was uniformly catastrophic ($R^2 = -2.00$).
This failure is biologically grounded: gliomas are infiltrative and
irregular, while meningiomas are well-circumscribed and dural-based. Shape
semantics differ qualitatively between the two tumour types, making a shared
shape representation a category error.

\begin{table}[H]
  \centering
  \caption{Shape probe $R^2$ across experimental conditions.}
  \label{tab:shape_failure}
  \begin{tabular}{lccl}
    \toprule
    \textbf{Condition} & \textbf{Linear $R^2$} & \textbf{RBF $R^2$} &
      \textbf{Cross-domain} \\
    \midrule
    Frozen baseline   & 0.08  & $\approx 0$ & $-2.00$ \\
    MEN LoRA $r$=8    & 0.11  & $\approx 0$ & $-2.00$ \\
    Dual LoRA $r$=8   & $-0.06$ & $\approx 0.15$ & $-2.00$ \\
    \bottomrule
  \end{tabular}
\end{table}

\textbf{Decision:} Drop the shape partition entirely from the SDP.

\subsubsection{Location: Temporally Static}

Meningiomas are extra-axial, dural-attached tumours that do not migrate. The
centroid coordinates at $t_0$ are essentially identical at subsequent
timepoints (modulo asymmetric growth). Modelling location as a time-varying
GP dimension fits a constant function, which is statistically vacuous and
wastes degrees of freedom. Within-domain location $R^2$ is reasonable
(0.35 MEN, 0.68 GLI), confirming the information is present but does not
need temporal modelling.

\textbf{Decision:} Relocate location from temporal latent dimension to static
GP covariate, entering the mean function as
$m_i(t) = \beta_0 + \beta_1 t + \boldsymbol{\gamma}^\top \mathbf{c}_i$,
where $\mathbf{c}_i$ is the centroid from patient $i$'s initial scan.

\subsubsection{Volume Target Simplification}

The 4-dimensional volume target $[\log V_{\text{total}}, \log V_{\text{NCR}},
\log V_{\text{ED}}, \log V_{\text{ET}}]$ is problematic because: (i)~label
semantics differ across tumour types (GLI oedema $\neq$ MEN oedema);
(ii)~sub-compartment segmentation quality is poor (TC Dice 0.40 on MEN);
(iii)~whole-tumour volume is the only clinically relevant growth endpoint for
meningiomas.

\textbf{Decision:} Reduce to a single target
$v^* = \log(V_{\text{WT}} + 1) \in \R^1$.

\subsubsection{Revised Architecture Summary}

\begin{equation}
  \text{R1:}\quad \mathbf{z} \in \R^{128} =
    [\underbrace{\Vol \in \R^{32}}_{\text{supervised}}
    \mid \underbrace{\Res \in \R^{96}}_{\text{unsupervised}}],
  \quad \pi_{\text{vol}}: \R^{32} \to \R^1.
  \label{eq:r1_partition}
\end{equation}

The PA-MOGP simplifies to a standard MOGP on $\R^{32}$ with a single
Mat\'{e}rn-5/2 kernel and ICM coregionalisation matrix
$\mathbf{B} = \mathbf{W}\mathbf{W}^\top + \diag(\boldsymbol{\kappa})$.

% ---------------------------------------------------------------------------
% 9.5 LoRA Marginality Discussion
% ---------------------------------------------------------------------------
\subsection{LoRA Marginality}
\label{sec:eval:marginality}

A notable finding is the marginal benefit of LoRA adaptation for downstream
feature quality. The frozen BrainSegFounder baseline achieves
$R^2 = 0.794$ for volume probing, while the best LoRA condition adds only
$\Delta R^2 = +0.016$ (within noise). This suggests that BrainSegFounder's
SSL pretraining on 41,400+ subjects already produces features with strong
meningioma volume decodability.

Several hypotheses explain this observation:
\begin{enumerate}[leftmargin=*]
  \item \textbf{SSL pretraining captures domain-general features.} The UK
    Biobank pretraining stage exposes the encoder to massive anatomical
    diversity, producing features that generalise across brain pathologies.
  \item \textbf{LoRA's low-rank constraint limits adaptation.} With
    $r = 8$ and $<0.32\%$ trainable parameters, LoRA may not have
    sufficient capacity to substantially restructure the 768-dimensional
    feature space.
  \item \textbf{Segmentation is a weak proxy for semantic quality.}
    Improving segmentation Dice does not necessarily improve the linear
    decodability of aggregate semantic properties (volume, location).
\end{enumerate}

This finding has important practical implications: for meningioma growth
prediction, the frozen BrainSegFounder encoder may be a sufficient starting
point, with the SDP and GP models providing the primary value.
